% =============================================================================
% Decoupling acceptance from scaling: an eta sweep at fixed eta1.
%
% Standalone. Own preamble, no bibliography, no \cite. The macro block below
% mirrors the one near the top of _Thesis_FINAL/Article3/Survey-part3-v1.tex,
% through \providecommand so that this file can later be folded into that
% document without a clash.
% =============================================================================
\documentclass[11pt,a4paper]{article}

\usepackage[T1]{fontenc}
\usepackage[utf8]{inputenc}
\usepackage[british]{babel}
\usepackage{amsmath, amssymb, amsthm}
\usepackage{booktabs}
\usepackage{graphicx}
\usepackage{geometry}
\usepackage{float}
\usepackage{enumitem}
\usepackage[hidelinks]{hyperref}
\geometry{margin=2.6cm}

% ---- macros mirrored from Survey-part3-v1.tex -------------------------------
\providecommand{\defeq}{\mathrel{:=}}
\providecommand{\kbar}{\overline{\kappa}}
\providecommand{\calO}{{\mathcal{O}}}
\providecommand{\calS}{{\mathcal{S}}}
\providecommand{\calU}{{\mathcal{U}}}
\providecommand{\NN}{{\mathbb{N}}}
\providecommand{\abs}[1]{\left \vert #1 \right \vert}
\providecommand{\norm}[1]{\left \lVert #1 \right \rVert}
\providecommand*\mykappa[1]{ \kappa_{\scriptscriptstyle \text{#1}} }
\providecommand*\myeta[1]{ \eta_{\scriptscriptstyle \text{#1}} }

\providecommand{\Rdelta}{\textsf{R-delta}}
\providecommand{\Rstep}{\textsf{R-step}}
\providecommand{\Rdfo}{\textsf{R-DFO}}
\providecommand{\Rgrad}{\textsf{R-grad}}
\providecommand{\Rrtr}{\textsf{R-RTR}}
\providecommand{\Rgrtr}{\textsf{R-gRTR}}
\providecommand{\RAdaptiveGrad}{\textsf{R-adaptive-grad}}
\providecommand{\RAdaptstep}{\textsf{R-adaptive-step}}
\providecommand{\Rdeltastep}{\textsf{R-delta-step}}
\providecommand{\Rtruncgrad}{\textsf{R-trunc-grad}}

\providecommand{\gone}{\gamma_1}
\providecommand{\gtwo}{\gamma_2}
\providecommand{\gthree}{\gamma_3}

\newtheorem{remark}{Remark}

% The conditions of Part~I carry names rather than numbers. \label inside an
% enumitem \item[] would pick up an unrelated counter, so the reference target
% is written explicitly.
\makeatletter
\newcommand{\namedlabel}[2]{%
  \protected@write\@auxout{}{\string\newlabel{#1}{{#2}{\thepage}{#2}{#1}{}}}%
}
\makeatother

\title{Decoupling acceptance from scaling\\
       {\large what it costs to accept a step the radius does not trust}}
\author{}
\date{}

\begin{document}
\maketitle

\begin{abstract}
The three thresholds of a trust-region method answer different questions. The
acceptance threshold $\eta$ decides whether a step is taken. The scaling
thresholds $\eta_1$ and $\eta_2$ decide whether the region shrinks or grows.
Setting $\eta < \eta_1$ opens a band on which a step is accepted and the radius
still contracts. We sweep $\eta$ from $0$ to $\eta_1$ for eight radius rules on
185 unconstrained CUTEst problems, at two values of $\eta_1$, and measure the
occupancy of that band before measuring anything else. PLACEHOLDER-ABSTRACT
\end{abstract}

% =============================================================================
\section{What Part~I says about \texorpdfstring{$\eta = 0$}{eta = 0}}
\label{sec: part I}
% =============================================================================

The question decides what the rest of this report means. Either the structural
conditions on $\{\Delta_k\}$ supply the quantified decrease that the classical
complexity argument takes from $\eta$, or the $\calO(\varepsilon^{-2})$ bound
does not survive $\eta = 0$. We read Part~I to settle it. Every label below is
quoted from that source.

Part~I fixes $0 \leq \eta \leq \eta_1 \leq \eta_2 < 1$ and defines
$\calS = \{k \geq 0 : \rho_k \geq \eta\}$ and
$\calU = \{k \geq 0 : \rho_k < \eta\}$. The abstract states the extension
directly. Part~I revisits the framework of Curtis and Scheinberg and of Wang and
Yuan, and extends it to three distinct scaling factors with decoupled step
acceptance, covering $\eta = 0$. The body attributes the Lyapunov framework to
the survey of Curtis and Scheinberg as developed by Cartis and Scheinberg and by
Blanchet, Cartis, Menickelly and Scheinberg, and records that the extension
allows $\eta = 0$, which is not treated there.

Part~I also states the classical position plainly. Several authors have used
$\eta = 0$, accepting any step that decreases the objective. That choice ensures
only that at least one accumulation point is first-order critical, which is
$\liminf_{k \to \infty}\norm{g_k} = 0$. Part~I cites Yuan's one-dimensional
example, which cycles among three points, two of them non-stationary.

\paragraph{The decrease is quantified by a different threshold.}
Three conditions carry the analysis. Condition~\ref{cond: LbR} bounds the radius
from below. Condition~\ref{cond: DecR} forces contraction after a poor step.
Condition~\ref{cond: ExpDecR} caps expansion after a good one.

\begin{description}[leftmargin=2.2em]
  \item[\textnormal{Condition C.LbR}]\namedlabel{cond: LbR}{C.LbR}
        \texttt{cond: delta\_k lower bound}. There is
        $0 < \mykappa{lbd} \leq 1$ with
        $\Delta_k \geq \frac{\mykappa{lbd}}{M_k}\min_{0 \leq i \leq k}\norm{g_i}$
        for all $k$, where $\{M_k\}$ is non-decreasing and bounds the model
        Hessian norm.
  \item[\textnormal{Condition C.DecR}]\namedlabel{cond: DecR}{C.DecR}
        \texttt{cond: decrease of the trust-region radius}. There are
        $\underline{\gtwo} \in (0,1)$ and $\myeta{dec} > 0$ such that
        $\rho_k < \myeta{dec}$ implies
        $\Delta_{k+1} \leq \underline{\gtwo}\,\Delta_k$.
  \item[\textnormal{Condition C.Exp/DecR}]\namedlabel{cond: ExpDecR}{C.Exp/DecR}
        \texttt{cond: increase of the trust-region radius}. There are
        $\underline{\gtwo} \in (0,1)$, $\overline{\gthree} > 1$ and
        $\myeta{inc} > 0$ such that $\rho_k < \myeta{inc}$ implies
        $\Delta_{k+1} \leq \underline{\gtwo}\,\Delta_k$, and
        $\rho_k \geq \myeta{inc}$ implies
        $\Delta_{k+1} \leq \overline{\gthree}\,\Delta_k$.
\end{description}

The thresholds $\myeta{dec}$ and $\myeta{inc}$ belong to the radius rule. They
are required to be strictly positive by the conditions themselves. They are not
$\eta$.

This is the whole answer. Part~I's liminf result,
\texttt{thm: Liminf of the gradient norm under strong assumption}, assumes
Condition~\ref{cond: LbR} and Condition~\ref{cond: DecR} with
$\myeta{dec} \geq \eta$ and $\myeta{dec} \neq 0$, and concludes
$\liminf_{k \to \infty} \norm{g_k} = 0$. At $\eta = 0$ the requirement
$\myeta{dec} \geq \eta$ is satisfied by every admissible rule, and
$\myeta{dec} \neq 0$ is a property of the rule rather than of the acceptance
test. The liminf result therefore holds at $\eta = 0$. The companion result
under the weaker Hessian assumption,
\texttt{thm: Liminf of the gradient norm under weaker assumption}, is stated the
same way with $\myeta{inc} \geq \eta$ and $\myeta{inc} \neq 0$.

The complexity bound is stated in $\myeta{dec}$ as well. Part~I's
\texttt{thm: Upper bound successful iter BTR} gives
\begin{equation}
  \abs{\{k < T_\varepsilon \, : \, \rho_k \geq \myeta{dec}\}}
  \;\leq\;
  \frac{2\big(f(x_0) - \mykappa{lbf}\big)\big(L + \mykappa{umh}\big)}
       {\myeta{dec}\,\mykappa{mdc}\,\mykappa{lbd}}\,\varepsilon^{-2} .
  \label{eqn: complexity}
\end{equation}
The denominator carries $\myeta{dec}$, not $\eta$. The classical argument turns
a successful iteration into a decrease of at least $\eta$ times the predicted
reduction and is vacuous at $\eta = 0$. Part~I's argument turns an iteration
with $\rho_k \geq \myeta{dec}$ into a decrease of at least $\myeta{dec}$ times
the predicted reduction. Such an iteration is accepted whenever
$\myeta{dec} \geq \eta$, so the decrease is realised. The route is the
structural conditions, and \eqref{eqn: complexity} is not vacuous at $\eta = 0$.

\paragraph{One discrepancy, which we report rather than resolve.}
The hypotheses of \texttt{thm: Upper bound successful iter BTR} open with
Algorithm~1 applied with $\eta > 0$. The bound it proves and the proof it gives
use $\eta$ only through $\myeta{dec} \geq \eta$, which is vacuous at $\eta = 0$.
Two readings are available. Under the first, the theorem is stated for
$\eta > 0$ and says nothing at $\eta = 0$. Under the second, the hypothesis
$\eta > 0$ is redundant and the conclusion holds whenever
$\myeta{dec} > 0$. Choosing between them changes what Part~I claims, so we
choose neither. The author should settle it.

\paragraph{What $\eta = 0$ does cost.}
The upgrade from $\liminf$ to $\lim$ is lost, and Part~I says so. Its
\texttt{th: firt-order TR CV} assumes $\eta > 0$ and concludes that
$\liminf_{k \to \infty}\norm{g_k} = 0$ implies
$\lim_{k \to \infty}\norm{g_k} = 0$. That theorem holds independently of the
radius rule, so no condition on $\{\Delta_k\}$ recovers it. The remark that
follows records the failure at $\eta = 0$. For iterations with
$\rho_k \in (\eta, \eta_1)$ the algorithm can accept steps with arbitrarily
small decrease, which admits accumulation points with non-zero gradient norm.

\paragraph{Consequence for this report.}
The band is a design freedom Part~I permits and the framework of Curtis and
Scheinberg forbids. It is not a practical gain bought against a lost complexity
bound. It is a practical question asked inside a theory that already covers it,
and the one guarantee at stake is the convergence of the whole sequence rather
than of a subsequence. Every measurement below should be read against that.

% =============================================================================
\section{What the solver does}
\label{sec: solver}
% =============================================================================

A rejected step leaves the iterate in place. The next iteration reuses the
gradient and pays one objective evaluation. An accepted step moves the iterate
and pays a gradient as well. Lowering $\eta$ converts rejected iterations into
accepted ones, so it converts cheap iterations into expensive ones.

We confirmed this from the source rather than assuming it. In
\texttt{src/Trust-region/common.jl} the acceptance test is
\texttt{st.accepted = st.$\rho$ >= p.$\eta$}, and the call
\texttt{grad!(nlp, c.x, c.g)} sits inside the branch \texttt{if st.accepted}.
The model update \texttt{update\_model!} and the curvature estimate are inside
the same branch. Outside the iteration the gradient is evaluated once, in
\texttt{\_init\_state!}. The only other gradient call is in
\texttt{\_refresh\_state!}, which the two sampled solvers use after a resample
and which no run in this report reaches.

The gradient is therefore evaluated once per accepted iteration and never on a
rejected one. An iteration count is not a work count here, and the two can move
in opposite directions. We report function, gradient and Hessian-vector
evaluations separately for that reason, and label iteration counts secondary
wherever they appear.

% =============================================================================
\section{Design}
\label{sec: design}
% =============================================================================

\paragraph{The sweep.}
For each rule we hold everything else fixed and sweep $\eta$ over
$\{0,\ \tfrac14\eta_1,\ \tfrac12\eta_1,\ \tfrac34\eta_1,\ \eta_1\}$. The last
is the classical coupling and is the baseline of every comparison.

Two strata. The first takes $\eta_1 = 0.1$, the value the campaign of Part~III
uses. The second takes $\eta_1 = 0.4$, because a band of width $0.1$ in $\rho$
may be too narrow to separate anything. Changing $\eta_1$ changes the rule's own
behaviour, so the second stratum is a separate axis. We report the two
separately and never pool them.

\paragraph{The rules.}
Part~III carries two rosters. Its appendix table lists eleven radius rules by
what the radius is tied to, which ratio selects the branch, whether the factor
is fixed or adaptive, and where each is first stated. Its body table lists
eleven and marks eight of them tested. The eight are $\Rdelta$, $\Rstep$,
$\Rdfo(\omega)$, $\Rgrad(\omega, \overline{\mu})$, $\Rgrad(\omega)$, $\Rrtr$,
$\RAdaptstep$ and $\RAdaptiveGrad(\omega)$. Those eight carry the primary
tables. They are the package's \texttt{RULES}, held at one set of constants,
$\gone = 0.25$, $\gtwo = 0.5$, $\gthree = 2.0$.

The two rosters disagree on their eleventh entry, and the disagreement bears on
the count, so we record it. The body table names $\Rdeltastep$,
$\Rgrtr(\omega,\overline{\mu})$ and $\Rtruncgrad$ as the three untested rules.
The appendix table names $\Rdeltastep$, $\Rgrtr(\omega,\overline{\mu})$ and
$\RAdaptiveGrad(\omega,\overline{\mu})$, and carries no entry for
$\Rtruncgrad$. Searching the source, $\Rtruncgrad$ has a macro in the preamble
and one appearance in the body table, and no environment defines it anywhere.
The author's own note in the source flags the same gap. We follow the appendix
table, which defines every rule it names.

Three rules are therefore surveyed and untested, and all three are implemented.
They are $\Rdeltastep$, $\RAdaptiveGrad(\omega,\overline{\mu})$ and
$\Rgrtr(\omega,\overline{\mu})$, carried by the package types
\texttt{RDeltaStep}, \texttt{RAdaptiveGradCapped} and \texttt{RRTRGrad}. We
report them in Section~\ref{sec: untested} and never inside the primary tables.

No package rule sits outside the surveyed taxonomy. Eleven rules are surveyed
by the appendix table and eleven package types implement them. The second-order
wrapper and its five aliases are not separate rules. They supply the criticality
measure $\omega$ and forward the update unchanged.

\paragraph{Constructor refusals.}
\texttt{validate\_thresholds} refuses $\eta_1 = 0$ for the step-driven rules
$\Rstep$, $\RAdaptstep$ and $\Rdeltastep$. This sweep varies $\eta$ and holds
$\eta_1$ at $0.1$ or $0.4$, so $\eta_1 = 0$ never arises and no refusal was
expected. None occurred, and the archive records the empty list.

\paragraph{Counting the band.}
Occupancy is counted on accepted iterations as $\rho_k < \eta_1$, from the
solver's own trace,
\begin{verbatim}
    count(i -> accepted[i] && rho[i] < eta1, eachindex(accepted))
\end{verbatim}
We do not count radius decreases. $\Rdelta$ contracts by $\gtwo$ on every mildly
successful iteration, so a shrinking radius does not isolate the band. The count
is zero by construction whenever $\eta = \eta_1$, and we report that column as
the check that the instrument works.

\paragraph{The problem set.}
CUTEst, unconstrained, $2 \leq n \leq 200$, which realises 185 problems. This is
the set Section~5 of Part~III reports. The package's \texttt{default\_problems()}
falls back to eight analytic problems when the CUTEst query returns empty, which
would produce a table that reads as a CUTEst benchmark without being one. The
experiment file checks \texttt{HAS\_CUTEST} and the realised count and stops
rather than falling back. The realised list of names and dimensions is written
to \texttt{tables/exp14\_problem\_set.txt} beside the results.

The model Hessian is exact and the subproblem solver is truncated conjugate
gradient throughout. \texttt{ExactMS} is admissible on all 185 problems at
$n \leq 200$, and we do not use it, so that no subsolver switch is confounded
with the swept threshold.

\paragraph{Pairing.}
The comparison is $\eta$ against $\eta_1$ within a rule on the same problem, so
it is paired. A pair enters the log ratio only when both settings solve the
problem. Problems solved by exactly one setting are named individually instead,
because a log ratio over survivors would drop exactly the cases that matter.
The tie band is $\abs{\log(w_\eta / w_{\eta_1})} \leq 0.05$.

PLACEHOLDER-RESULTS

\end{document}
